\subsection{Data Splitting, Validation, and Benchmarking}

Evaluation estimates how the trained system behaves on data not used to fit it.

\begin{definition} \label{def:sl-train-validation-test} \hlt{Training, Validation, and Test Splits}\\
The training set fits parameters, the validation set selects models and hyperparameters, and the test set estimates final performance after selection.  Mixing these roles produces optimistic error estimates.
\end{definition}

\begin{definition} \label{def:sl-benchmark} \hlt{Benchmark}\\
A benchmark specifies data, preprocessing rules, train-test splits, metrics, and allowed external resources.  Without these conventions, reported performance numbers need not be comparable.
\end{definition}

\begin{remark} \label{rem:sl-data-leakage} \hlt{Data Leakage}\\
Leakage occurs when information from validation or test data enters training, feature construction, preprocessing, hyperparameter selection, or manual iteration.  Leakage is a systems failure as much as a statistical one.
\end{remark}

\begin{remark} \label{rem:sl-metric-choice} \hlt{Metric Choice}\\
Accuracy, log loss, calibration error, ranking metrics, cost-sensitive utility, and latency measure different aspects of a system.  A benchmark metric should match the decision problem it stands in for.
\end{remark}
